\section{Methodology}

Our methodology is designed to analyze TTC reliability as a spatiotemporal pattern rather than a single KPI. The workflow integrates historical delay logs, GTFS reference data, and GTFS-Realtime service alerts into a layered analytics pipeline and a stakeholder-facing dashboard. The end-to-end architecture is shown in Figure~\ref{fig:ttc-architecture}.

\subsection{Approach Overview}
The overall approach has six stages:
\begin{enumerate}
    \item Ingest historical bus/subway delay data, static GTFS, and live GTFS-RT service alerts.
    \item Standardize schemas and canonicalize noisy route/station identifiers.
    \item Build mode-specific analytical units (route-first for bus, station-first for subway).
    \item Aggregate disruptions across time windows (year/month/weekday/hour bins).
    \item Compute explainable reliability metrics and composite rankings.
    \item Validate historical hotspots using live GTFS-RT alert alignment.
\end{enumerate}

\begin{figure}[t]
\centering
\begin{tikzpicture}[
    node distance=0.45cm and 0.65cm,
    box/.style={
        rectangle,
        rounded corners,
        draw=black,
        thick,
        align=center,
        minimum width=2.2cm,
        minimum height=0.8cm,
        fill=gray!5
    },
    arrow/.style={
        -{Latex[length=2mm]},
        thick
    }
]

\node[box] (hist) {Historical\\Delay Logs};
\node[box, right=of hist] (gtfs) {Static\\GTFS};
\node[box, right=of gtfs] (alerts) {GTFS-RT\\Service Alerts};

\node[box, below=of gtfs] (clean) {Clean \&\\Normalize};
\node[box, below=of clean] (agg) {Space-Time\\Aggregation};
\node[box, below left=of agg] (metrics) {Reliability\\Metrics};
\node[box, below right=of agg] (validation) {Live Alert\\Validation};

\node[box, below=1.8cm of agg] (dash) {Linked Dashboard\\(Overview, Hotspots, Time, Causes, Drill-Down, Live)};

\draw[arrow] (hist) -- (clean);
\draw[arrow] (gtfs) -- (clean);
\draw[arrow] (alerts) -- (clean);
\draw[arrow] (clean) -- (agg);
\draw[arrow] (agg) -- (metrics);
\draw[arrow] (agg) -- (validation);
\draw[arrow] (metrics) -- (dash);
\draw[arrow] (validation) -- (dash);

\end{tikzpicture}
\caption{TTC Pulse data-analytics architecture: ingestion, canonicalization, aggregation, reliability scoring, and live-alert validation feeding the dashboard.}
\label{fig:ttc-architecture}
\end{figure}

\subsection{Data-Analytics Methodology}
\textbf{Data preparation and normalization.}
Historical delay events are transformed into a canonical event structure with service date/time, mode, route or station context, incident category, and operational measures (\texttt{Min Delay}, \texttt{Min Gap}). GTFS dimensions are used as a reference layer to standardize route/station identity and preserve consistency across analytical views.

\textbf{Spatiotemporal analysis design.}
We aggregate by entity and time to answer ``where'', ``when'', and ``how often'' disruptions recur. Temporal views include yearly, monthly, weekday, and hour-bin slices. This is appropriate for TTC reliability because recurrence patterns are strongly schedule- and period-dependent.

\textbf{Reliability metric framework.}
For each entity $u$ in time window $t$, we compute:
\begin{itemize}
    \item \textbf{Frequency} ($F$): total disruption-event volume in the entity-time slice (implemented as summed \texttt{event\_count}/\texttt{frequency}); here, an \emph{incident} means one recorded delay/disruption event tied to a service date and incident category
    \item \textbf{Severity} ($\mathrm{Sev}_{90}$): upper-tail delay behavior (90th-percentile delay minutes)
    \item \textbf{Regularity} ($\mathrm{Reg}_{90}$): upper-tail gap/headway irregularity (90th-percentile gap minutes)
    \item \textbf{Cause Mix} ($C$): disruption cause composition signal
\end{itemize}

Cause mix is computed from category concentration:
\[
C(u,t)=1-\sum_{k} p_k(u,t)^2,
\]
where $p_k$ is the share of category $k$ within entity-time slice $(u,t)$.

For ranking, we use the implemented composite formula:
\[
S(u,t)=0.35\,z(F)+0.30\,z(\mathrm{Sev}_{90})+0.20\,z(\mathrm{Reg}_{90})+0.15\,C.
\]
The z-score normalization for each metric $x$ is:
\[
z(x)=\frac{x-\mu_x}{\sigma_x},
\]
computed within the peer comparison set (same mode/date-window context). If $\sigma_x=0$, the implementation safely falls back to zero contribution for that standardized term. The dashboard always exposes component metrics beside the composite score.

Interpretation of $S(u,t)$ is relative (not bounded to $[-1,1]$):
\begin{itemize}
    \item $S(u,t)>0$: worse-than-peer reliability profile in the selected comparison set
    \item $S(u,t)<0$: better-than-peer reliability profile in the selected comparison set
    \item $S(u,t)\approx 0$: near peer-average reliability profile
\end{itemize}
Larger absolute magnitude means stronger deviation from peer-average behavior.

\subsection{Data-Analytics Architecture}
The implementation follows a local-first layered architecture:
\begin{itemize}
    \item \textbf{Raw/Bronze}: ingestion and lineage-preserving storage
    \item \textbf{Silver}: canonical dimensions, bridges, review queues, normalized facts
    \item \textbf{Gold}: stakeholder marts for reliability metrics, rankings, temporal profiles, and alert validation
\end{itemize}

Runtime stack:
\begin{itemize}
    \item \textbf{Storage/Query}: DuckDB + Parquet
    \item \textbf{Serving}: Streamlit dashboard
    \item \textbf{Live capture}: GTFS-RT service-alert sidecar scheduled every 30 minutes (local OS scheduler)
\end{itemize}

\textbf{Why we shifted from Airflow to local OS scheduling.}
Based on course-scope guidance discussed with the instructor (keep deployment reproducible on student laptops and avoid server-grade operational dependencies), we moved from Airflow-first orchestration to local OS-native scheduling (\texttt{launchd} on macOS and Task Scheduler on Windows). The key reasons were:
\begin{itemize}
    \item only one recurring production task was needed (30-minute alert sidecar cycle)
    \item Airflow introduces extra runtime services (scheduler/webserver/metadata DB) that are disproportionate for this scope
    \item local schedulers improve teammate portability and reduce setup/maintenance overhead
    \item this choice preserves focus on analytics quality rather than infrastructure operations
\end{itemize}
Airflow DAGs are retained as legacy/reference artifacts, but active execution for this project milestone is local-first.
\textbf{Scope expansion note:} streetcar support has started at the data-model level and will be integrated into reliability marts and dashboard views in the next phase.

\subsection{Algorithms and Implemented Components}
\textbf{(1) Canonicalization strategy.}
Bus analytics use route-first canonicalization and subway analytics use station-first canonicalization. Ambiguous/unmatched mappings are retained in review-oriented tables instead of being silently dropped.

\textbf{(2) Windowed ranking aggregation.}
For selected service-date windows, rankings are recomputed from metrics marts using windowed aggregation and rank ordering on composite score (with frequency as tie-break context).

\textbf{(3) Alert-hotspot alignment metric design (planned KPI, partial implementation).}
Let $A$ be the set of live alerts in a validation window and let $H_K$ be the historical top-$K$ hotspots from the selected baseline window. The evaluation metrics are:
\[
\mathrm{HitRate}=\frac{1}{|A|}\sum_{a \in A}\mathbf{1}[a \in H_K],
\qquad
\mathrm{Lift}=\frac{\mathrm{HitRate}}{\mathrm{BaselineHitRate}}.
\]
where \(\mathrm{BaselineHitRate}\) is computed against a non-hotspot baseline set of equal size (or a random baseline under the same selector constraints).

Interpretation for stakeholders:
\begin{itemize}
    \item \textbf{HitRate}: ``Of all current live alerts, what fraction fall in historically high-risk hotspots?''
    \item \textbf{Lift}: ``How much better is hotspot targeting than a baseline expectation?''
    \item \(\mathrm{Lift}>1\): live alerts are more concentrated in historical hotspots than baseline (supports recurrence claim)
    \item \(\mathrm{Lift}\approx 1\): hotspot concentration is similar to baseline (weak validation)
    \item \(\mathrm{Lift}<1\): historical hotspots do not explain current live alerts well in that window
\end{itemize}
This directly supports the TTC Pulse storyline by testing whether ``today's disruptions'' align with ``historically recurring hotspots''.

Current implementation status at midterm: the dashboard currently surfaces alert-match status and selector-validity diagnostics in the Live Alert Alignment page, while explicit \(\mathrm{HitRate}\)/\(\mathrm{Lift}\) KPIs are specified and ready for the evaluation stage.

\subsection{Limitations}
\begin{itemize}
    \item Cause-mix signal is less stable in very fine-grain slices.
    \item GTFS-RT selector sparsity can reduce exact entity-match confidence.
    \item Spatial certainty is stronger for canonical subway stations than for bus route-centroid approximations.
    \item Composite ranking interpretation requires component-metric visibility to avoid over-reliance on a single score.
\end{itemize}
